\chapter{Maintenance failure: the slow accident}
\label{vol10:ch14}

\epigraph{The accident did not happen on a single day; it happened over the decades of deferral that preceded the day.}{}

\chapmeta{Half-life: durable. Archetypes: failure, ethics. Exercise target: 15. Project track: hybrid.}

\noindent
Maintenance failure is the slow accident: the gradual erosion of margins, capabilities, and physical condition over years and decades until the system can no longer perform what its design specified, and the failure that physics or chemistry was patiently producing finally occurs. The mechanisms are usually not exotic (corrosion, fatigue, wear, scale, fouling, drift); the conditions that allow them to accumulate are the political and economic and organisational conditions that defer the maintenance the system requires. This chapter is about that pattern: how maintenance failures develop, what makes the developed-economy infrastructure stock so vulnerable to them, and how the cases (Genoa Morandi, Flint, ageing nuclear) read.

\section{Maintenance neglect as a leading failure cause}\pathtag{core}

The accident-investigation record consistently identifies maintenance failures among the leading proximate causes of major failures. The Bhopal disaster (Chapter 10, \cite{acc:icmr-bhopal-2004}) is partly a maintenance failure: the refrigeration system that would have held the MIC tanks at $0\,\degC$ had been disabled to save operating cost; the flare tower was not operational; the vent gas scrubber was out of service. The TMI-2 accident (Chapter 9, \cite{acc:kemeny-tmi-1979}) includes a maintenance failure as its initiating mechanism: the auxiliary feedwater discharge valves had been left closed after maintenance. The BP Texas City refinery explosion \cite{acc:csb-bp-texas-city-2007} included multiple maintenance shortfalls in the affected unit's instrumentation and protection systems. The Boeing 787 lithium-ion battery incidents \cite{acc:ntsb-aar14-01-787-battery} had maintenance dimensions in the airline-level response to early incidents. The Aloha 243 incident \cite{acc:ntsb-aar94-04} was substantially a maintenance failure of the fleet inspection regime in addressing the multi-site fatigue cracking it had documented but not addressed.

The pattern is that maintenance neglect contributes to a substantial fraction of major engineering accidents, but the contribution is rarely visible at the moment of the accident. The accident reports identify the immediate technical mechanism (the relief valve that stuck open; the joint that failed; the line that ruptured), and the institutional analyses identify the upstream conditions (the production pressure that deferred the inspection; the budget cut that reduced the maintenance workforce; the organisational structure that did not surface the maintenance backlog to decision-makers). The accident's date is when the physics finally exceeded the system's degraded capability; the failure was years or decades in preparation.

The structural reasons that maintenance is systematically under-attended:

\textbf{Maintenance is invisible when it works.} A well-maintained system operates without incident; the absence of incident is not credited to the maintenance budget; the budget that produces nothing visible is the budget that is cut first. The unrewarded-success pattern is characteristic of all preventive activity (safety; security; preventive medicine) and applies to maintenance with particular force.

\textbf{Maintenance has long time constants.} The deferral of a maintenance action in a given budget cycle does not produce visible consequences in the same budget cycle; the consequences accumulate over years. The decision-maker who defers the action does not bear the consequence; the decision-maker who eventually receives the failure inherits a problem decades in preparation.

\textbf{Maintenance is mundane.} The professional incentives in engineering organisations reward new design, novel solutions, visible projects; they do not strongly reward the routine work of keeping existing systems operating. The career path of an engineer who specialises in maintenance is typically less prestigious and less remunerative than the career path of an engineer who designs new systems; the talent allocation reflects the incentive structure.

\textbf{Maintenance budgets are politically vulnerable.} Public infrastructure maintenance is funded by tax revenue or by user fees; both are politically contested. The political and budget cycles are shorter than the maintenance time constants; the decision-makers who set the budgets are not the decision-makers who will face the failures the budget cuts produce.

\textbf{The conditions for maintenance neglect compound.} An organisation that defers maintenance reduces the workforce capable of performing it; the reduced workforce produces backlog; the backlog increases the unit cost of catch-up; the catch-up is again deferred; the system condition deteriorates further. The dynamic is self-reinforcing in the direction of deterioration.

The engineering response is to make the maintenance state visible (instrumentation; condition monitoring; reporting); to make the maintenance economics legible (lifecycle cost analysis; deferred-maintenance accounting; failure-probability-as-a-function-of-deferral models); and to defend the maintenance budget structurally (statutory minimums; protected funding streams; independent assessment of maintenance status). The institutional record is mixed; the pattern of maintenance failure recurs because the structural conditions that produce it have not been reliably defeated.

\begin{principle}[Maintenance is a system property]
A system's reliability is not the property of its as-built design but the property of its as-maintained condition. The design specifies the upper bound; the maintenance regime determines how close to the upper bound the system actually operates. An engineering organisation that does not include the maintenance regime in its system reliability analysis is analysing a system different from the one it actually operates.
\end{principle}

\section{Aging infrastructure}\pathtag{core}

Substantial portions of the developed-economy infrastructure stock were built in the 1950s through 1970s, with design lives of 50 to 100 years and with maintenance assumptions that have not been consistently fulfilled. The stock is reaching ages at which design life is exhausted or substantial intervention is required to extend it.

The US bridge stock provides a representative dataset. The American Society of Civil Engineers (ASCE) Infrastructure Report Card (the most recent edition is 2021) grades US infrastructure across multiple categories on a scale of A through F; the 2021 overall grade for US infrastructure is C-, an improvement from 2017's D+. Bridges specifically receive a C grade in 2021; the report identifies approximately 7.5\% of US bridges (43,578 of approximately 617,000) as "structurally deficient" (a categorisation indicating that one or more major structural components is in poor condition and warrants near-term action). The deficiency does not mean the bridge is unsafe (structurally deficient bridges remain in service under reduced load ratings, with monitoring, or with planned replacement schedules), but it does indicate that the stock requires substantial investment to remain in service.

The US water-infrastructure stock has similar dynamics. The 2021 ASCE report card grades drinking water C-, wastewater D+, and stormwater D. The American Water Works Association and other industry analyses identify approximately a trillion-dollar US water-infrastructure investment requirement over the next 25 years to replace aging mains, treatment plants, and distribution systems. The Flint water crisis (Chapter 14 case below) is one expression of the broader pattern; the lead-service-line inventory in older US cities is a known vulnerability whose remediation cost is substantial and whose remediation timeline depends on political and budgetary decisions still unfolding.

The European infrastructure stock has parallel dynamics. The post-war reconstruction (1950s-1970s) built much of the modern European infrastructure; the design lives are approaching their nominal limits; the maintenance investment required to extend the stock or to replace it is the subject of ongoing political negotiation. The Genoa Morandi Bridge \cite{acc:mit-morandi-2018} is an instance of the broader pattern in the Italian motorway network; similar dynamics exist in other European national networks.

The Asian infrastructure stocks are at different points in their lifecycles. The Japanese stock includes substantial post-war construction that is reaching maintenance-intensive ages; the rapid Chinese infrastructure expansion of the 1990s-2010s has produced a younger stock with maintenance issues that are smaller in current intensity but will scale substantially over coming decades.

The cumulative pattern is that maintenance of ageing infrastructure is one of the dominant engineering challenges of the current and coming decades, larger in aggregate engineering work than the construction of new infrastructure. The discipline of the chapter is to make the work as visible as new construction, to recognise the failures (Genoa Morandi; Flint; the I-35W Mississippi River bridge collapse in 2007; multiple dam failures of the 2020s) as evidence of the broader maintenance shortfall, and to argue for the institutional structures that would address it.

\section{Deferred-maintenance dynamics}\pathtag{standard}

The dynamics of deferred maintenance can be modelled with simple structural equations that reveal the failure modes.

Let $M(t)$ denote the maintenance backlog at time $t$, $S(t)$ the system's effective condition (1 = as-new; 0 = unserviceable), $b(t)$ the per-period maintenance budget, $c(t)$ the cost of the maintenance work the system requires per period to hold condition constant. The standard accounting identity is:

\[
\frac{dM}{dt} = c(t) - b(t)
\]

and the system condition evolves as:

\[
\frac{dS}{dt} = -\alpha (1 - S) - \beta M
\]

where $\alpha$ is the rate of condition deterioration from age alone and $\beta$ is the rate at which backlog itself accelerates deterioration (deferred work compounds: a leaking joint that is not repaired now produces additional damage; a corroded connection that is not addressed produces additional corrosion). The model is linear and simplified; richer models add stochastic shocks (storm events; load excursions; equipment failures) and nonlinear thresholds (the system's condition crossing the level at which catastrophic failure becomes substantially more probable).

Under steady-state operation with $b(t) = c(t)$, $M$ is constant and $S$ approaches a steady-state value determined by $\alpha$ and the initial conditions. Under budget compression $b(t) < c(t)$, $M$ grows linearly, $S$ falls; the system's failure probability rises; eventually the budget that would catch up requires substantial increase, the political and financial barriers to which produce the structural difficulty of recovering from a deferred-maintenance pattern.

The dynamic produces several characteristic failure modes:

\textbf{Backlog blow-up.} The backlog grows over years to a magnitude that no single budget cycle can address; the catch-up requirement becomes structurally impossible without substantial new revenue. The catch-up cost is typically several times the cost of the routine maintenance that would have prevented the backlog accumulation; the political acceptance of the catch-up is harder than the original routine investment would have been.

\textbf{Capability erosion.} The reduced maintenance budget reduces the workforce capable of performing maintenance; the trained technicians leave or are not replaced; the institutional knowledge erodes; the productivity per dollar of maintenance budget falls; the same dollar buys less maintenance than it did. The dynamic is self-reinforcing: a workforce-reduced maintenance organisation cannot effectively absorb a budget increase even if one is provided, because the capability to use the increase has been lost.

\textbf{Cliff failure.} The system condition crosses the threshold at which catastrophic failure probability rises sharply; the failure occurs; the political and financial consequences of the failure are vastly greater than the maintenance budget that would have prevented it. The Genoa Morandi failure is the canonical recent example: the catch-up investment in the post-failure replacement bridge plus the legal, regulatory, and political costs vastly exceeded what continuous adequate maintenance would have cost.

\textbf{Capability migration.} The deferred-maintenance pattern in one part of the infrastructure stock produces failures that political response addresses; the response captures attention and budget; other parts of the stock continue to defer maintenance with no political attention; the next failure occurs in a different part of the stock. The pattern produces a sequence of "surprises" that are not actually surprising to anyone who has been tracking the underlying maintenance state.

The defences are: explicit maintenance accounting (the maintenance state is tracked, reported, and audited; the backlog magnitude is visible); statutory funding minimums (the maintenance budget cannot be cut below a specified floor without explicit override); life-cycle cost analysis (the budget decision is presented in terms of total ownership cost over the asset's lifecycle, including the failure cost from inadequate maintenance, not only the current-period maintenance expense); and independent assessment (the maintenance state is reviewed by parties not subject to the same budget pressures as the operating organisation).

\section{The political economy of maintenance}\pathtag{standard}

Maintenance is a political and economic activity as much as a technical one; the engineering analysis of why maintenance is under-attended is incomplete without the political economy that produces the under-attention.

The political-economic dynamics:

\textbf{Visibility asymmetry.} New construction is visible (a new bridge; a new building; a new highway); maintenance is invisible (the bridge that does not fall down; the building that does not leak; the highway that does not pothole). Political incentives favour the visible; the politician who cuts a ribbon on a new project receives credit; the politician who funds maintenance receives no comparable credit unless the maintenance prevents a visible failure (which is by definition invisible). The asymmetry biases political decision-making toward new construction at the expense of maintenance of the existing stock.

\textbf{Cost-shifting across electoral cycles.} The decision-maker who cuts maintenance in the current budget cycle does not bear the consequence in the current cycle. The consequence accumulates over years; the decision-maker may have left office before the consequence is felt; the political incentive to defer maintenance is concentrated in the current cycle while the consequence is dispersed across future cycles. The dynamic is the standard one of inter-temporal political-economic decision-making, with maintenance being an especially clean example.

\textbf{Privatisation and concession dynamics.} When infrastructure is operated under concession or private ownership, the operator's incentive depends on the concession structure. A concession with no end-of-concession condition requirements gives the operator the incentive to minimise maintenance toward the end of the concession period; a concession with strict end-condition requirements gives the operator the incentive to invest sufficiently. The Italian motorway concession structure under which the Morandi Bridge was operated had been the subject of contested analysis on these points before the bridge collapse \cite{acc:mit-morandi-2018}; the post-collapse regulatory response included revision of the concession framework.

\textbf{Public-budget structural compression.} Public budgets are dominated by mandatory expenditure (pensions, healthcare, debt service); the discretionary portion in which infrastructure maintenance competes for funds is a small fraction of total budgets. The political tendency is to cut the discretionary portion under fiscal pressure; infrastructure maintenance shares the cut with other discretionary categories.

\textbf{User-fee dynamics.} Infrastructure funded by user fees (water utilities; road tolls; transit fares) faces political resistance to fee increases. The fee structures often do not generate enough revenue to fund the asset's lifecycle maintenance; the resulting funding shortfall accumulates as deferred maintenance.

\textbf{Federalism and intergovernmental funding.} In federal systems, infrastructure funding is distributed across multiple levels of government; the funding responsibility for maintenance may be unclear or contested; the level of government that owns the asset may not be the level that funds the maintenance. The funding shortfalls and their consequences are dispersed across the federal-state-local structure, making accountability diffuse.

The engineering profession's role in the political economy of maintenance is multiple: provide the technical analysis that makes the maintenance state legible to non-technical decision-makers; advocate for the institutional structures (statutory minimums; independent assessment; lifecycle cost analysis) that protect maintenance funding; and participate in the post-failure investigations that establish the connection between deferred maintenance and the failures it produces. The work is at the intersection of engineering and public-policy disciplines; engineers who confine themselves to the technical work do not address the political-economic conditions that produce the failures they investigate.

\section{Case: Genoa Morandi Bridge, Flint water crisis, ageing nuclear plants}\pathtag{core}

\subsection*{Genoa Morandi Bridge, 14 August 2018}

The Polcevera Viaduct in Genoa, Italy, collapsed at 11:36 on 14 August 2018, killing 43 people \cite{acc:mit-morandi-2018}. The bridge was a hybrid cable-stayed design from 1967 in which the main stays consisted of post-tensioned concrete-encased steel tendons; the design was unusual and had become harder to inspect and maintain as the bridge aged. The southern stays of Pylon 9 had experienced progressive corrosion in the bridge's marine and industrial atmosphere; the corrosion was difficult to detect by the non-destructive testing technologies available at inspection and was the subject of contested inspection findings and contested recommendations for intervention. On 14 August, the deteriorated stays ruptured; the load redistribution exceeded the remaining stays' capacity; the pylon and 210 m of deck collapsed.

The maintenance-failure analysis identifies several patterns. The bridge had been the subject of inspections that identified deterioration of the stays; the inspection findings had not led to an action plan on a timescale that would have prevented the failure. The contracted inspection consultants had recommended more aggressive intervention; the concession operator's internal decision-making had not endorsed the recommendations. The regulator had not compelled the operator to act on the inspection findings. The bridge's age (51 years at collapse) was within its nominal design life, but the maintenance investment required to keep it operating safely had not been made.

The post-collapse response illustrates the cost asymmetry. The replacement Genoa San Giorgio Bridge (designed by Renzo Piano; opened 2020) cost approximately EUR~200 million. The economic disruption to the Port of Genoa and to the broader Ligurian economy during the bridge's outage was substantial. The legal and regulatory consequences (criminal prosecutions; concession revocation; framework renegotiation) imposed additional costs. The cumulative cost to the Italian state, the city, the operator, the affected businesses, and the victims' families substantially exceeded what continuous adequate maintenance would have cost across the bridge's operating life. The case is the textbook example of cliff failure: the system condition crossed the threshold; the failure cost was vastly greater than the prevention cost; the political and budgetary attention that the failure compelled was attention that earlier funding would not have received.

The institutional response in Italy includes the revision of the motorway-concession framework, the renationalisation of major segments of the network, the development of more aggressive inspection regimes for cable-stayed bridges of the Morandi generation, and the broader political recognition that the Italian infrastructure stock requires substantial maintenance investment that prior decades have under-funded.

\subsection*{Flint water crisis, 2014-2019}

The Flint water supply was switched from Detroit-system water to Flint River water on 25 April 2014; the absence of corrosion-control treatment caused lead from the city's lead service lines to leach into drinking water; approximately 100,000 residents were exposed to elevated lead levels for approximately 18 months before the contamination was acknowledged and the supply was switched back \cite{acc:fwatf-2016}. The Flint Water Advisory Task Force identified the technical mechanism (chloride-to-sulfate-ratio-driven corrosion of the lead service lines in the absence of orthophosphate treatment), the organisational mechanism (a state-appointed emergency manager whose financial-recovery priorities drove the supply switch; MDEQ regulatory advice that corrosion-control treatment was not required, an advice inconsistent with EPA's Lead and Copper Rule), and the public-health consequences (elevated blood-lead levels in children; a Legionnaires' disease outbreak that killed at least 12 people).

The maintenance-failure framing applies along multiple dimensions. The city's lead service-line inventory was itself the result of historical maintenance and replacement decisions; the lead service lines were no longer installed after the 1980s, but the existing inventory had not been replaced because replacement was expensive and the operating water-treatment regime had been adequate to keep lead levels within action thresholds. The water-treatment-plant capacity to treat Flint River water adequately would have required investment that the city, under state-appointed emergency management focused on financial recovery, did not make. The state regulator's capacity to oversee the technical complexity of the supply transition was limited by long-term underfunding of the MDEQ drinking-water program. The cumulative pattern is maintenance failure at multiple levels: the service-line inventory; the treatment-plant capability; the regulator's oversight capability.

The post-crisis response illustrates the cost asymmetry as starkly as the Morandi case. Lead service-line replacement costs in Flint, the long-term health monitoring of affected residents, the legal settlements, the criminal prosecutions, the political and reputational costs to the affected officials and institutions, the substantial federal and state intervention costs all aggregate to a sum vastly larger than what continuous corrosion-control treatment and proactive service-line replacement would have cost over the same period. The federal Lead and Copper Rule revisions (final 2021) include strengthened requirements that respond to the Flint mechanism; the broader US response to lead service lines has accelerated through the 2021 Infrastructure Investment and Jobs Act funding.

\subsection*{Ageing nuclear plants}

The US commercial nuclear power fleet comprises plants whose construction dates fall mostly in the 1970s and 1980s. The original licences were 40 years; the NRC's licence renewal process allows extension to 60 years and, more recently, to 80 years. The plants that have operated for several decades are subject to ageing-management programs that explicitly address the mechanisms (radiation embrittlement of reactor vessels; thermal and irradiation aging of cable insulation; concrete degradation in containment structures; fatigue cycling of primary coolant piping) that would otherwise produce the kinds of failures the discipline of this chapter is about.

The institutional discipline that has developed is substantial. The NRC's Generic Aging Lessons Learned (GALL) report catalogues the ageing mechanisms and the inspection and intervention regimes that address them. The INPO peer-review process examines plant maintenance and ageing-management practices. The licence renewal process requires plants to demonstrate that their ageing-management programs are adequate to the extended operating period. The institutional discipline is the response to TMI \cite{acc:kemeny-tmi-1979} and to the broader recognition that the original construction's design life was conservative but the maintenance investment required to extend operations safely was substantial.

The cases in which the discipline has failed are also instructive. The Davis-Besse reactor's vessel-head corrosion incident of 2002 (in which boric-acid corrosion had eaten through approximately six inches of the carbon-steel reactor vessel head, leaving only a quarter-inch of stainless-steel cladding holding back primary coolant pressure) is an example of an ageing mechanism (boric-acid corrosion at small primary-system leak points) that the established inspection regime had not caught. The post-incident regulatory response addressed the inspection regime gap; the broader lesson was that ageing-management programs require continued evolution as the operating experience accumulates and new mechanisms surface.

The ageing-nuclear-plant case is the positive counterpoint to Genoa Morandi and Flint: an industry with substantial regulatory oversight, with an established peer-review discipline, with a maintenance investment that the regulatory structure compels and the institutional incentives support, operating ageing infrastructure with substantially better-managed risk than the comparable civil-infrastructure stock. The lesson is partly that the institutional conditions for adequate maintenance can be established (the nuclear industry's post-TMI evolution is the canonical example) and partly that the conditions are not the default; without active institutional construction, the default is the Morandi pattern.

\begin{principle}[Inspection is not maintenance]
A bridge that is inspected and identified as deteriorating is not a bridge that has been maintained. The inspection produces information; the maintenance is the action on the information. An institutional structure in which inspection findings reliably produce action is the structure that prevents Morandi-class failures; an institutional structure in which inspection findings can be deferred indefinitely is the structure that produces them. The engineering profession's contribution is to design and defend the institutional structures that close the inspection-to-action loop.
\end{principle}

\begin{archetype}[failure]
Maintenance failure is the slow-time expression of the failure archetype. The conditions for failure develop over years; the institutional response to the developing conditions is the determinant of whether the failure becomes catastrophic. The accident's date is when the physics finally exceeds the deteriorated capability; the failure was years in preparation.
\end{archetype}

\begin{project}[Maintenance assessment of one piece of regional infrastructure]
Select one piece of public infrastructure in the reader's region (a bridge; a water-treatment plant; a power substation; a section of motorway or rail line; a public building) about which the reader can locate publicly available inspection or condition-assessment information. Conduct a maintenance-status assessment following the discipline of this chapter: identify the asset's age, design specifications, and operating history to the extent the public record allows; locate the most recent inspection reports; identify the inspection findings and the institutional response; estimate the maintenance backlog (the work the inspections identify that has not been done); identify the funding source for the asset's maintenance and the political-economic dynamics that shape it; identify what a maintenance-failure-class accident in this asset would cost (lives, property, economic disruption); compare to the cost of addressing the identified backlog. Deliverable: a four-page document with the assessment, an opinion on whether the asset's current maintenance trajectory is adequate, and (if it is not) a specific recommendation for action.

\textbf{Hazard.} The exercise may reveal that the asset's condition is more concerning than the reader assumed; the political and reputational consequences of publishing such an assessment are real. The professional register of the assessment matters; advocacy without analytical depth is not engineering work and is properly identified as advocacy.
\end{project}

\section{Exercises}

\subsection*{Calculation}\pathtag{core}

\begin{exercise}
For the deferred-maintenance dynamic model with $\frac{dM}{dt} = c - b$ and $\frac{dS}{dt} = -\alpha(1-S) - \beta M$, compute the steady-state condition $S^*$ when $b < c$ (constant budget shortfall). Identify the dependence of $S^*$ on $\alpha$, $\beta$, and the budget shortfall $c - b$.
\end{exercise}

\begin{exercise}
The Genoa San Giorgio Bridge replacement cost approximately EUR~200 million. If continuous adequate maintenance of the original Morandi Bridge would have cost an estimated EUR~5 million per year, compute the number of years of maintenance the replacement cost equates to. Compare to the bridge's 51-year service life at collapse.
\end{exercise}

\begin{exercise}
The chloride-to-sulfate mass ratio (CSMR) in Flint River water was approximately 1.6; in Detroit-system water it was approximately 0.4. Industry guidance suggests CSMR above approximately 0.5 is associated with elevated lead release in lead service lines under non-corrosion-controlled conditions. Compute the ratio of the two CSMRs. Discuss the role of corrosion-control treatment in keeping lead release low even at elevated CSMR.
\end{exercise}

\subsection*{Derivation}\pathtag{standard}

\begin{exercise}
Derive a model of the capability-erosion dynamic: a maintenance workforce of size $W$ produces $\eta W$ units of maintenance per period; the workforce attrites at rate $\delta$ and grows by recruitment at rate $r(b)$ where $b$ is the maintenance budget. Identify the conditions under which a budget reduction produces irreversible workforce loss and the consequent productivity-per-dollar reduction.
\end{exercise}

\begin{exercise}
Derive the expected lifecycle cost of an asset under two maintenance scenarios: (a) continuous adequate maintenance with cost $c$ per year and no end-of-life intervention; (b) reduced maintenance with cost $b < c$ per year and end-of-life replacement at cost $R$. Identify the conditions on $R$, $c$, $b$, and the asset's life under each scenario at which scenario (a) is preferred.
\end{exercise}

\subsection*{Estimation}\pathtag{standard}

\begin{exercise}
Estimate the aggregate US infrastructure maintenance backlog (the catch-up cost to bring infrastructure to a "state of good repair") using the most recent ASCE Infrastructure Report Card. Compare to current annual federal infrastructure spending.
\end{exercise}

\begin{exercise}
Estimate the cost of lead service-line replacement in Flint per service line; multiply by the estimated inventory of US lead service lines (approximately 9 million per recent estimates). Compare to the annual federal water-infrastructure spending.
\end{exercise}

\subsection*{Simulation}\pathtag{standard}

\begin{exercise}
Implement the deferred-maintenance dynamic model. Sweep budget shortfall and stochastic shock magnitude; identify the regimes in which the system condition reaches the cliff-failure threshold within 10 years, 30 years, and 50 years.
\end{exercise}

\begin{exercise}
Simulate a population of 1000 infrastructure assets with independent age and condition trajectories under a fixed total maintenance budget. Identify the allocation rule (worst-condition-first; statutory-priority; oldest-first; risk-weighted) that maximises the population-average condition over a 50-year horizon.
\end{exercise}

\subsection*{Design}\pathtag{standard}

\begin{exercise}
Design a maintenance-state accounting framework for a municipal water utility. Specify the condition metrics, the reporting cadence, the public disclosure requirements, and the funding-decision integration. Justify the design against the Flint case \cite{acc:fwatf-2016}.
\end{exercise}

\begin{exercise}
Design a statutory minimum-maintenance-budget mechanism for a national motorway system. Specify the floor (as a function of asset value or replacement cost), the override conditions, the audit and enforcement mechanisms, and the political-economic dynamics the mechanism would change.
\end{exercise}

\subsection*{Judgement}\pathtag{mastery}

\begin{exercise}
Argue: the Genoa Morandi Bridge collapse \cite{acc:mit-morandi-2018} was not a maintenance failure but a design failure (the bridge's hybrid cable-stayed design was inherently maintenance-impractical, and no realistic maintenance regime could have prevented the failure). Argue the opposite. Identify what the dispute implies about the boundary between design and maintenance failures.
\end{exercise}

\begin{exercise}
Discuss the role of state-appointed emergency management in the Flint water crisis \cite{acc:fwatf-2016}. Identify the structural conditions under which financial-recovery prioritisation produces technical risk; identify the institutional changes that would protect against the dynamic.
\end{exercise}

\begin{exercise}
Argue: the nuclear industry's ageing-management discipline is transferable to civil infrastructure and would substantially reduce the rate of Morandi-class failures. Identify the obstacles to the transfer (institutional, financial, political) and the conditions under which the transfer could occur.
\end{exercise}

\begin{exercise}
A municipal infrastructure authority is asked to prioritise its limited maintenance budget across a portfolio of ageing bridges. Develop a defensible prioritisation framework. Identify the engineering inputs, the political-economic constraints, and the ethical considerations.
\end{exercise}
